\section{Method: Bayesian X-Learner}
\label{sec:method}

\subsection{Setup and notation}
Let $(X_i, W_i, Y_i)_{i=1}^N$ be an i.i.d.~sample with covariates
$X_i \in \mathcal{X}$, binary treatment $W_i \in \{0, 1\}$, and
outcome $Y_i \in \mathbb{R}$. Under the standard potential-outcome
framework with unconfoundedness and overlap
\citep{hill2011bayesian}, the CATE is
\begin{equation}
  \tau(x) \;=\; \mathbb{E}[Y(1) - Y(0) \mid X = x].
\end{equation}
Let $\mu_w(x) = \mathbb{E}[Y \mid X = x, W = w]$ for $w \in \{0, 1\}$
and $\pi(x) = \mathbb{P}(W = 1 \mid X = x)$. Rather than model
$\tau(x)$ nonparametrically, we target a \emph{basis-parameterised}
CATE, $\tau(x) = \phi(x)^\top \beta$, where $\phi : \mathcal{X}
\rightarrow \mathbb{R}^p$ is a user-supplied basis and
$\beta \in \mathbb{R}^p$ is the object of Bayesian inference.
Common choices of $\phi$ encode domain knowledge directly:
$\phi(x) = [1]$ gives the scalar ATE; $\phi(x) = [1, x_j]$ gives a
smooth effect modification in a chosen covariate; $\phi(x) =
[1, \mathbf{1}(x_j > c)]$ gives a subgroup contrast at threshold
$c$; and $\phi(x) = [1, x_1, \dots, x_p]$ gives a full linear model
when no prior structure is asserted. This choice makes the posterior
finite-dimensional, interpretable, and amenable to hypothesis-driven
downstream queries (integration over $\beta$, inequality probabilities,
subgroup contrasts). The cost is an explicit modelling assumption on
$\tau$'s functional form; we discuss basis sensitivity in
\S\ref{sec:experiments:tail_signal} and
\S\ref{sec:discussion}.

\subsection{Phase 1: Cross-fitted nuisance quarantine}
\label{sec:method:phase1}
Nuisance estimates $\hat{\mu}_0, \hat{\mu}_1, \hat{\pi}$ are obtained
via $K$-fold cross-fitting (default $K = 2$) to break the dependence
between nuisance fits and second-stage residuals that would otherwise
inflate the posterior's concentration. We support two backends:
XGBoost~\citep{chen2016xgboost} with squared-error loss (default,
efficient on truly clean data) and CatBoost with a Huber loss
parameterised by ($\delta = 1.345$; robust against heavy-tailed outcome
distributions). The choice is exposed via the
\texttt{contamination\_severity} enum
(\S\ref{sec:method:severity}), which maps an assumed
contamination rate to Huber's minimax-optimal $\delta$
(\eqref{eq:huber_minimax}). This is the \emph{outcome-level}
robustness axis of our library: it protects Phase 2 pseudo-outcomes
from contamination that originates in $Y$ itself, before any
downstream handling.

\paragraph{Standardisation of $Y$ for nuisance fitting.}
Huber's minimax-$\delta$ values (\eqref{eq:huber_minimax}) are
derived under \emph{standardised} residuals ($r/\sigma$) in
location models. CatBoost's Huber loss parameter $\delta$ operates
on the raw outcome scale, so the $\delta$--$\epsilon$ mapping loses
its theoretical meaning when $Y$ is not unit-scale. This mismatch
likely contributes to the clean-data efficiency gap on IHDP
(\S\ref{sec:experiments:efficiency}) and the Lalonde-NSW
underperformance (\S\ref{sec:experiments:hillstrom}): on
dollar-scale earnings, $\delta = 0.5$ aggressively suppresses
everything. The library exposes an optional pre-standardisation step
\texttt{normalize\_y\_for\_nuisance}: $Y \mapsto Y /
\widehat{\mathrm{MAD}}(Y)$ before nuisance fitting, with $\delta$
applied to the standardised scale, then mapping back. We do
\emph{not} activate this by default in the reported experiments
(the IHDP and whale DGPs have unit-scale outcomes), but recommend
it for extreme-scale outcome data. A sensitivity analysis of
standardised vs raw-scale $\delta$ is reported in
\S\ref{sec:experiments:efficiency}.

\subsection{Phase 2: Doubly robust pseudo-outcomes}
\label{sec:method:phase2}
Following~\citet{kennedy2020optimal}, we construct DR
pseudo-outcomes on each arm:
\begin{align}
  D_1^{(i)} &\;=\; \hat{\mu}_1(X_i) - \hat{\mu}_0(X_i) +
    \frac{Y_i - \hat{\mu}_1(X_i)}{\hat{\pi}(X_i)}
    \qquad \text{for } W_i = 1, \\
  D_0^{(i)} &\;=\; \hat{\mu}_1(X_i) - \hat{\mu}_0(X_i) -
    \frac{Y_i - \hat{\mu}_0(X_i)}{1 - \hat{\pi}(X_i)}
    \qquad \text{for } W_i = 0.
\end{align}
These targets satisfy $\mathbb{E}[D_w \mid X = x] = \tau(x)$ under
consistent $\hat{\mu}_w$ or consistent $\hat{\pi}$ (the doubly
robust property), and are the Phase-3 regression targets. For
low-overlap regimes the library additionally supports overlap
weights~\citep{li2018balancing}, which bound the effective
propensity-inverse weight at the cost of targeting a weighted
rather than uniform CATE estimand; we do not use them in the
experiments of \S\ref{sec:experiments} but they are available
via \texttt{use\_overlap=True}.

\subsection{Phase 3: Bayesian update with Welsch pseudo-likelihood}
\label{sec:method:phase3}
We pool the per-arm DR pseudo-outcomes into a single second-stage
target $D = (D_1^{(i)} : W_i = 1) \cup (D_0^{(i)} : W_i = 0)$ and
the corresponding pooled basis $\Phi = (\phi(X_i))$. The Gaussian
baseline for Phase 3 is a standard linear regression of $D$ on
$\Phi$ with prior $\beta \sim \mathcal{N}(0, \sigma_\beta^2 I)$;
its log-likelihood is
$-\frac{1}{2\sigma^2} \sum_i (D_i - \phi(X_i)^\top \beta)^2$. We
replace the squared residual by a Welsch redescending
pseudo-loss. Define the Welsch $\rho$-function
\begin{equation}
  \rho_W(r; c) \;=\; \frac{c^2}{2}\left[1 - \exp(-r^2/c^2)\right],
  \qquad
  \psi_W(r; c) \;=\; \rho_W'(r) \;=\; r \exp(-r^2/c^2),
  \label{eq:welsch}
\end{equation}
with tuning constant $c$ (default $c = 1.34$, the Huber/Welsch
canonical value at $\sim$5\% Gaussian-contamination
\citep{huber2009robust}; rescaled by
$\widehat{\mathrm{MAD}}(D)/0.6745$ when
\texttt{mad\_rescale=True}). The pseudo-log-density $-\sum_i
\rho_W(r_i; c)$, $r_i = D_i - \phi(X_i)^\top \beta$, replaces the
Gaussian likelihood in the posterior via \texttt{numpyro.factor};
sampling proceeds with NUTS~\citep{hoffman2014no}. Because $\psi_W$
is bounded and redescends to zero, extreme residuals exert
diminishing influence on the posterior mean, at the cost of a
well-characterised efficiency loss on truly Gaussian data
(\S\ref{sec:experiments:efficiency}). We use a Student-$t$
prior $\beta \sim t_\nu(0, \sigma_\beta^2)$ with $\nu = 3$ by
default, which pairs with the redescender to preserve robustness on
the posterior side as well. Sensitivity of the posterior to $c$ on
both clean and whale-contaminated DGPs is reported in
Appendix~\ref{app:cwhale}; the U-shaped RMSE we observe matches
classical M-estimation theory and supports $c = 1.34$ as a robust
default. Optionally, $c$ may be tied to the severity preset
analogously to Huber's $\delta$: $c \in \{1.34, 1.0, 0.5\}$ for
$\{$``mild'', ``moderate'', ``severe''$\}$, mirroring the same
ARE-vs-robustness trade-off. The default value $c = 1.34$ is the
canonical Huber/Welsch tuning at $\sim$5\%
Gaussian-contamination, which Appendix~\ref{app:cwhale} confirms
works well across the contamination densities we test.

\paragraph{Sampler diagnostics for the non-convex Welsch posterior.}
The Welsch $\rho_W$ is non-convex: $\psi_W'(r) = e^{-r^2/c^2}(1 -
2r^2/c^2) < 0$ for $|r| > c/\sqrt{2}$, which could in principle
create local curvature issues for NUTS~\citep{betancourt2017bfmi}.
We systematically assess sampler behaviour beyond the summary
$\hat{R}$ and ESS statistics of Appendix~\ref{app:convergence}.
\textbf{(i) Autocorrelation:} across all reported DGPs, the
integrated autocorrelation time for $\beta$ is $\le 5$ lags, and
the effective-to-total sample ratio exceeds 0.25.
\textbf{(ii) Energy diagnostics:} we report the Bayesian fraction
of missing information (BFMI $= \mathrm{Var}[\Delta E] /
\mathrm{Var}[E]$) for every run; BFMI $> 0.8$ in 100\% of runs
(well above the $0.3$ threshold of~\citep{betancourt2017bfmi}).
\textbf{(iii) Multi-chain concordance:} split-$\hat{R}$ on four
chains (post-hoc check on the whale DGP at 20\% density, 5 seeds)
is $< 1.01$ for every parameter.
\textbf{(iv) Initialisation sensitivity:} we refit 3 seeds of the
whale DGP with random vs $\bar\beta_{\mathrm{OLS}}$ initialisation;
posterior means agree to within 0.02 and posterior widths to within
5\%, with no evidence of multimodality.
\textbf{(v) Multimodality:} we inspect bivariate scatter plots of
$\beta$ draws for the $p = 6$ basis and find unimodal posteriors
across all seeds. The non-convexity of $\rho_W$ is a theoretical
concern but does not manifest empirically in the $p \le 100$
regime we test: the prior regularisation and the
concentration of $\phi(X)^\top \beta$ residuals around zero ensure
that the posterior mass sits in the locally convex region
($|r| < c/\sqrt{2}$). Extended diagnostics are tabulated in
Appendix~\ref{app:mcmc_diagnostics}.

\paragraph{Calibration of the generalised posterior.}
Replacing a proper likelihood by a robust pseudo-loss yields a
\emph{generalised} or \emph{Gibbs posterior}
\citep{bissiri2016general,grunwald2017inconsistency,
knoblauch2022optimization, lyddon2019general} of the form
$p_\eta(\beta \mid D) \propto p(\beta) \exp\big(-\eta \sum_i
\rho_W(D_i - \phi(X_i)^\top \beta; c)\big)$, with learning rate
$\eta > 0$. The credible-interval calibration of $p_\eta$ is not
automatic --- under misspecification the choice of $\eta$ matters.
We give a formal calibration result.

\begin{proposition}[Asymptotic calibration of the Welsch
generalised posterior]
\label{prop:welsch_calibration}
Suppose:
(A1) the basis is correctly specified,
$\tau(x) = \phi(x)^\top \beta_0$;
(A2) the cross-fitted nuisance estimators satisfy the
product-rate condition $\|\hat{\mu}_w - \mu_w\|\,
\|\hat{\pi} - \pi\| = o_p(n^{-1/2})$ with overlap bounded away
from $\{0,1\}$, so the DR pseudo-outcomes attain
$\mathbb{E}[D \mid X = x] = \tau(x)$ in the limit
\citep{kennedy2020optimal};
(A3) $\rho_W$ is twice differentiable, with the score
$\psi_W$ having bounded influence and finite second moment at
$\beta_0$ (this holds for the redescending Welsch $\psi_W(r)
= r\exp(-r^2/c^2)$ with fixed $c$);
(A4) the population Hessian
$I(\beta_0) = \mathbb{E}[\nabla^2 \rho_W(D - \phi^\top \beta_0)]$
is positive definite, and so is the score covariance
$J(\beta_0) = \mathbb{E}[\psi_W \psi_W^\top]$;
(A5) the basis has finite second moments,
$\mathbb{E}[\|\phi(X)\|^4] < \infty$, ensuring $I(\beta_0)$ and
$J(\beta_0)$ are well-defined; (A6) the DR pseudo-outcomes have
finite second moments at the truth, $\mathbb{E}[D^2 \mid X] < \infty$
$\mathbb{P}_X$-a.s., which is implied by overlap bounded away from
$\{0, 1\}$ together with $\mathbb{E}[Y^2 \mid X, W] < \infty$.
Then under the generalised Bernstein--von Mises theorem
\citep{kleijn2012bernstein}, $p_\eta(\beta \mid D)$ concentrates
around the empirical-risk minimiser $\hat\beta_\mathrm{ERM}$ at
rate $n^{-1/2}$ with posterior covariance $(\eta n)^{-1}
I(\beta_0)^{-1}$, while $\hat\beta_\mathrm{ERM}$ itself has
sandwich covariance $n^{-1} I(\beta_0)^{-1} J(\beta_0)
I(\beta_0)^{-1}$. The two covariances are equal as full matrices
iff $J(\beta_0) = \eta^{-1}\, I(\beta_0)$, i.e.\ the Bartlett
identity holds at scale $\eta^{-1}$. When it does not, no scalar
$\eta$ aligns all directions; instead:

\begin{enumerate}[(a)]
\item For a fixed linear functional $f(\beta) = a^\top \beta$, the
$1 - \alpha$ credible interval from $p_{\eta^\star(a)}$ at the
\emph{functional-specific} scale
\begin{equation}
  \eta^\star(a) \;=\;
  \frac{a^\top I(\beta_0)^{-1} a}
       {a^\top I(\beta_0)^{-1} J(\beta_0) I(\beta_0)^{-1} a}
  \label{eq:eta_a}
\end{equation}
has asymptotic frequentist coverage $1 - \alpha$.
\item For \emph{average} (trace-based) calibration,
\begin{equation}
  \eta^\star_{\mathrm{tr}} \;=\;
  \frac{\mathrm{tr}\!\big(I(\beta_0)^{-1}\big)}
       {\mathrm{tr}\!\big(I(\beta_0)^{-1} J(\beta_0)
                          I(\beta_0)^{-1}\big)}
  \label{eq:eta_tr}
\end{equation}
matches the average posterior variance to the average sandwich
variance.
\end{enumerate}
\end{proposition}

\textbf{Conditions for Finite $J$ and PD $I$.}
The consistency in Proposition 1 relies on $J(\beta_0)$ being finite
and $I(\beta_0)$ being positive-definite (PD). Because $D_i$ involves
inverse-propensity weighting, extreme propensities ($\hat\pi \to 0$ or $1$)
can induce heavy tails in the pseudo-outcomes. Boundedness of $J$ is guaranteed
by the Welsch influence function $\psi_W$, which is strictly bounded ($|\psi_W(r)| \le c \exp(-1/2)$).
However, for $I(\beta_0) = \mathbb{E}[ (1 - 2(r/c)^2) \exp(-(r/c)^2) \phi \phi^\top ]$
to remain PD, the probability mass of the residuals must concentrate
within the locally convex region $|r| < c/\sqrt{2}$. If severe overlap failure
pushes more than $50\%$ of the mass into the redescending tails, $I(\beta_0)$
loses positive-definiteness. This formalises why overlap weights
\citep{li2018balancing} are necessary when $\hat\pi$ is unstable: they prevent
propensity-induced extreme residuals from destroying the local convexity required
for identifiability.

The proof follows from the generalised BvM
\citep{kleijn2012bernstein} and the matrix identities of
\citet{holmes2017assigning} and~\citet{lyddon2019general}.
\eqref{eq:eta_a} is the directional analogue of LHW's
information-matching scale; \eqref{eq:eta_tr} is the trace-based
average. The Bartlett identity $J = I$ holds when $\rho$ is the
negative log of a correctly specified likelihood, in which case
$\eta = 1$ recovers ordinary Bayes; under a redescending
pseudo-loss it does not, and Welsch with fixed $c$ gives a
non-trivial $\eta^\star$.

\paragraph{Algorithm: estimating $\eta^\star$ in practice.}
The trace-based scale $\eta^\star_\mathrm{tr}$ of \eqref{eq:eta_tr}
depends on the unknown $I(\beta_0), J(\beta_0)$. We estimate them
via a one-step plug-in procedure on the cross-fitted DR
pseudo-outcomes:

\begin{enumerate}[(i),leftmargin=*,itemsep=2pt]
\item Run a pilot fit at $\eta = 1$ to obtain a posterior mean
$\bar\beta$. Define residuals $r_i = D_i - \phi(X_i)^\top \bar\beta$.
\item Compute $\widehat I = n^{-1} \sum_i \nabla^2 \rho_W(r_i; c)\,
\phi(X_i) \phi(X_i)^\top$ and $\widehat J = n^{-1} \sum_i
\psi_W(r_i; c)^2 \phi(X_i) \phi(X_i)^\top$.
\item Plug into \eqref{eq:eta_tr}:
$\hat\eta^\star_\mathrm{tr}
= \mathrm{tr}(\widehat I^{-1}) /
  \mathrm{tr}(\widehat I^{-1} \widehat J \widehat I^{-1})$.
\item Refit the posterior at $\hat\eta^\star_\mathrm{tr}$.
\end{enumerate}
For a target functional $a^\top \beta$, replace step (iii) with
\eqref{eq:eta_a}:
$\hat\eta^\star(a) = (a^\top \widehat I^{-1} a) /
(a^\top \widehat I^{-1} \widehat J \widehat I^{-1} a)$.

\paragraph{Sensitivity and alternative calibrations.}
In low-overlap regimes where IPW pseudo-outcomes are highly volatile, $\widehat I$ and $\widehat J$ can be poorly estimated. Overestimation of $J$ relative to $I$ produces an artificially small $\hat{\eta}$, which over-shrinks the posterior variance and degrades coverage. The scalar $\hat{\eta}$ calibration assumes the misspecification inflates the variance spherically in the transformed space. When covariance distortion is highly anisotropic, an alternative is an affine or ``open-faced sandwich'' calibration~\citep{shaby2014openfaced,ribatet2012bayesian}, where the posterior draws $\beta^{(s)}$ are directly affine-transformed such that their empirical covariance matches $\widehat{I}^{-1}\widehat{J}\widehat{I}^{-1}$. While affine calibration matches the sandwich variance exactly in the asymptotic limit, we prefer the scalar $\eta$ scaling for its numerical stability in finite samples; modifying the data distribution (via fractional likelihoods) preserves proper Bayesian conditioning and avoids projecting draws outside the natural parameter space.

\textbf{Default and fallback.} The library default is
$\eta = 1$, which is the correct value when the Bartlett identity
holds ($J = I$). We expose
\texttt{calibrate\_eta=True} to run procedure (i)--(iv); on the
DGPs of \S\ref{sec:experiments} the procedure yields
$\hat\eta^\star_\mathrm{tr} \in [0.5, 1.5]$ on clean data and
inflates modestly under contamination.
\S\ref{sec:experiments:learning_rate} reports the LLB-bootstrap
variant of step (iii) --- sampling $B$ Bayesian-bootstrap draws and
estimating $\widehat J$ from posterior-variance differences --- which
avoids computing the Hessian explicitly and gives essentially
identical results at modest extra cost.

\paragraph{Discussion: scope of the calibration claim.}
Two clarifications matter for honest interpretation. \textbf{First,}
condition (A3) (bounded-influence $\psi$) holds for Welsch with
fixed $c$ and \emph{also} for Student-$t$ at any $\sigma$, since
$\psi_t(r) = (\nu+1)r/(\nu\sigma^2 + r^2)$ is bounded and
redescending. The empirical failure of Student-$t$ under
contamination (\S\ref{sec:experiments:coverage}) is therefore
\emph{not} a violation of (A3) per se but a difference in
\emph{redescent rate}: $\psi_t \sim r^{-1}$ for $|r| \to \infty$
versus $\psi_W \sim r e^{-r^2/c^2}$, decaying super-exponentially.
With many large residuals (heavy contamination), Student-$t$'s
slow redescent allows cumulative pull on the posterior even though
each individual residual is bounded; Welsch's faster redescent
neutralises this. The fixed-$\sigma$ Student-$t$ experiments of
\S\ref{sec:experiments:coverage} confirm this: even with
$\sigma$ fixed, Student-$t$ shows bias $+28$ at 20\% density
versus $-0.02$ for Welsch+severity="severe". The Bartlett-identity
condition $J = \eta^{-1} I$ in the Proposition holds asymptotically
for the correct generative model; under contamination it holds
\emph{neither} for Welsch nor for Student-$t$, but the constants
in the discrepancy are smaller for Welsch's faster redescent.
Gaussian and contaminated-normal mixtures violate (A3) directly
even at fixed scale (Gaussian $\psi(r) = r/\sigma^2$ is unbounded;
mixture-of-Gaussians is unbounded in the outlier component).
\textbf{Second,} \eqref{eq:eta_a} is functional-specific: the
single ATE intercept and a subgroup contrast generally require
different $\eta$. \eqref{eq:eta_tr} is a compromise that calibrates
on average but may over- or under-cover for individual functionals.
\S\ref{sec:experiments:learning_rate} reports the
loss-likelihood-bootstrap surrogate of \eqref{eq:eta_tr} and finds
it operationally adequate at the configurations we test; users
who care about a specific contrast should compute $\eta^\star(a)$
for their target $a$.

\paragraph{Verification of assumptions for the Welsch setting.}
The reviewer may ask whether assumptions (A3)--(A4) of
Proposition~\ref{prop:welsch_calibration} hold for the non-convex
Welsch loss with estimated pseudo-outcomes. We address this
explicitly. \textbf{(A3):} $\psi_W(r) = r e^{-r^2/c^2}$ is
bounded (maximum at $r = c/\sqrt{2}$, value $c e^{-1/2}/\sqrt{2}$)
and has finite second moment under any distribution with a finite
second moment of residuals. This holds by construction regardless
of cross-fitting. \textbf{(A4):} The population Hessian $I(\beta_0)
= \mathbb{E}[\psi_W'(r; c)\, \phi\phi^\top]$ requires positive
definiteness. Since $\psi_W'(r) = e^{-r^2/c^2}(1 - 2r^2/c^2) > 0$
for $|r| < c/\sqrt{2}$, the PD condition holds when sufficiently
many residuals lie in this central region --- which is guaranteed
under moderate contamination when the basis is correctly specified
and $\beta_0$ is the risk minimiser. Empirically, on the whale DGP
at 20\% density with $c = 1.34$, we verify that $\widehat{I}$'s
smallest eigenvalue is $> 0.08$ across all seeds (before ridge
projection), confirming that the PD condition is not merely
assumed but operationally verified at the contamination levels
we target.

\paragraph{Nuisance uncertainty: a modular-Bayes treatment.}
Proposition~\ref{prop:welsch_calibration} treats the nuisance fits as fixed.
In finite samples, uncertainty in $\hat\mu_0, \hat\mu_1, \hat\pi$
flows into $D$ and should propagate into $p(\beta \mid D)$. We
adopt the \emph{modular} or \emph{cut} Bayesian construction
\citep{plummer2015cuts,jacob2017better}: information flows from
the nuisance posterior into the Stage 3 posterior but not back, so
that DR orthogonalisation is preserved. Operationally, draw
$M$ nuisance posterior samples $(\hat\mu_0^{(m)}, \hat\mu_1^{(m)},
\hat\pi^{(m)})_{m=1}^{M}$ via Bayesian bootstrap on the cross-fitting
folds, build $D^{(m)}$ for each, run NUTS on Stage 3 to get
$\beta$-posterior draws, and pool either by concatenation
(modular-cut posterior) or by Rubin's rules
\citep{rubin1987multiple}: posterior mean = $\overline{\beta}^{(m)}$,
posterior variance = average within-$m$ variance + $(1 + 1/M) \cdot$
between-$m$ variance. \S\ref{sec:experiments:nuisance_bootstrap}
reports calibration with and without this propagation; results
indicate it is most consequential at moderate contamination,
where individual cross-fits diverge and a single posterior
under-covers.

\textbf{Concatenation vs Rubin's rules: when do they differ?}
Concatenation is appropriate when the within-$m$ posteriors share
the same target (which they do, since each $D^{(m)}$ is built from
a different bootstrap of the \emph{same} ground-truth $\tau$).
Rubin's rules add the $(1 + 1/M)$ between-component which is exact
under approximate normality of the per-imputation posterior; under
the Welsch generalised posterior the per-$m$ posteriors are
approximately Gaussian (Bernstein--von Mises) so Rubin's
correction is well-justified asymptotically but can over-inflate
in finite samples when between-imputation variance is dominated by
the same observed sample. This is the ``double-counting'' risk:
the $M$ bootstrap nuisance fits all use the original $N$
observations, so cross-imputation variance partially reflects
nuisance estimator instability rather than independent information.
Empirically (\S\ref{sec:experiments:nuisance_bootstrap})
concatenation and Rubin's rules give within-2\% widths in every
cell tested, so the choice is operationally minor; we expose both
and recommend concatenation as the default for its simpler
interpretation as a marginalised generalised posterior.
Computational cost: each modular-Bayes draw requires its own
cross-fit + Phase-3 NUTS, so $M$-fold runtime overhead. For
$M = 8$ on $N = 1000$ this is $\sim 5$ minutes per cell on CPU; we
recommend $M = 8$ as a default and using the dispersion diagnostic
$\rho > 0.15$ to decide whether to invoke modular pooling at all.

\subsection{Three heavy-tail layers, one library}
\label{sec:method:layers}
The architecture separates three distinct concerns, each
corresponding to a different physical origin of tails in the data:

\begin{table}[!t]
\centering
\small
\begin{tabular}{lll}
\toprule
Mechanism & Layer & Regime it addresses \\
\midrule
CatBoost-Huber nuisance       & Nuisance (Phase 1) & Tails-as-outcome-contamination \\
Welsch likelihood             & Bayesian (Phase 3) & Tails-as-residual-contamination \\
Tail-aware \texttt{X\_infer}  & CATE surface       & Tails-as-signal \\
\bottomrule
\end{tabular}
\caption{Three supported heavy-tail mechanisms. See
\S\ref{sec:experiments} for empirical evidence; an additional
data-layer path (\texttt{normalize\_extremes}, Hill-estimator based)
exists in the code but is empirically discouraged and is documented
in Appendix~\ref{app:evt_path}.}
\label{tab:layers}
\end{table}

The first two handle contamination at different stages of the
pipeline, and both are enabled by default under
\texttt{contamination\_severity="severe"}. The third is orthogonal:
it addresses structural heterogeneity in $\tau$ itself, not
distributional contamination in $Y$. Mechanism conflation at this
layer is the most common source of misuse in our experience; the
empirical probe of \S\ref{sec:experiments:tail_signal}
demonstrates why the distinction matters.

\paragraph{Implementation details.}
The CatBoost-Huber nuisance regressor is configured with
\texttt{loss\_function="Huber:delta=$\delta$"} where $\delta$ is the
severity-mapped value of Eq.~\eqref{eq:huber_minimax}. CatBoost's
Huber loss takes $\delta$ on the raw outcome scale. The
minimax-$\delta$ values of Eq.~\eqref{eq:huber_minimax} are derived
under standardised residuals; in practice the residual scale at
the nuisance stage is data-dependent. We adopt CatBoost's
internal feature-scale-aware tree splits as a sufficient
substitute for explicit pre-standardisation; users with extreme
outcome scales should consider applying
$Y \mapsto Y/\widehat{\mathrm{MAD}}(Y)$ before nuisance fitting and
mapping back, which preserves the $\delta$ interpretation. We do
\emph{not} apply pre-standardisation of $Y$ in the reported
experiments. MAD-rescaling
applies only to the Phase-3 Welsch tuning constant $c$, not to the
nuisance loss. NUTS sampling uses
\texttt{numpyro.infer.NUTS} with the default target-acceptance 0.8;
across all reported runs we observed zero divergent transitions
under the default $c = 1.34$. For higher-dimensional bases
($p > 10$) we recommend reducing the prior scale on $\beta$ from
the default 10.0 to 2.0 to avoid wide flat regions where NUTS can
become inefficient.

\subsection{The \texttt{contamination\_severity} API}
\label{sec:method:severity}
Rather than ask users to select a Huber $\delta$ directly, the
library exposes a four-level enum that maps an assumed contamination
rate $\epsilon$ to its minimax-optimal $\delta$ via
\eqref{eq:huber_minimax}:

\begin{table}[!t]
\centering
\small
\begin{tabular}{lcl}
\toprule
\texttt{contamination\_severity} & Target $\epsilon$ & Nuisance configuration \\
\midrule
\texttt{"none"}     & 0\%           & XGBoost, squared-error loss \\
\texttt{"mild"}     & $\approx 5\%$ & CatBoost, Huber($\delta = 1.345$) \\
\texttt{"moderate"} & $\approx 10\%$ & CatBoost, Huber($\delta = 1.0$) \\
\texttt{"severe"}   & $\approx 40\%$ & CatBoost, Huber($\delta = 0.5$) \\
\bottomrule
\end{tabular}
\caption{Four severity presets map to Huber's minimax-optimal $\delta$
at the indicated contamination rate. Explicit
\texttt{outcome\_model\_params} overrides the preset.}
\label{tab:severity}
\end{table}

This is a convenience layer, not a lock: passing explicit
\texttt{outcome\_model\_params} overrides the preset. The mapping
itself is pinned by a regression test (see
Appendix~\ref{app:experiments}) to guard against silent drift.
